%% kapa: exam=1 round=34 year=2023 period=2 form=A subject=appraisal_law subject_ko=감정평가관계법규 seq=18 number=18
%% source: Q-Net 감정평가사 기출문제 (cst00309) artlSeq=5211789 file=raw/1008/5211789/2023년도 제34회 감정평가사 1차시험 2교시.pdf
%% conversion: kordoc 4.24.51 -> tools/gen1.py (auto); review=proofread-sonnet; flags=none
\begin{question}{18}
감정평가 및 감정평가사에 관한 법령상 과징금에 관한 설명이다. (　　)에 들어갈 내용을 순서대로 나열한 것은?
\begin{framedtext}
○ 감정평가법인에 대한 과징금부과처분의 경우 과징금최고액은 (　　)원이다. \par
○ 과징금납부의무자가 과징금을 분할납부하려면 납부기한 (　　)일 전까지 국토교통부장관에게 신청하여야 한다. \par
○ 과징금을 납부기한까지 납부하지 아니한 경우에는 납부기한의 다음 날부터 과징금을 (　　)까지의 기간에 대하여 가산금을 징수할 수 있다. \par
\end{framedtext}
\begin{choices}
  \item 3억, 10, 납부한 날의 전날
  \item 3억, 30, 납부한 날
  \item 5억, 10, 납부한 날의 전날
  \item 5억, 10, 납부한 날
  \item 5억, 30, 납부한 날의 전날
\end{choices}
\end{question}
\answer{3}
